\documentclass[11pt,a4paper]{article}
\usepackage[utf8]{inputenc}
\usepackage{geometry}
\geometry{margin=1in}
\usepackage{lineno}
\usepackage{xcolor}
\usepackage{titlesec}
\usepackage{array}

% Command for redacting personal information with a black box
\newcommand{\redacted}{\rule{2.5cm}{2ex}}

\title{\textbf{Social Science Fair Interview 04 \\ Transcript: Value and Pressure of Football Defenders (Central Midfielder Perspective)}}
\author{\textbf{Research Team Members:} Luke Sun, Kiara Wei, Winston Huang}
\date{May 28, 2026}

\begin{document}

\maketitle

\section*{Transcript Metadata}
\begin{tabular}{ll}
    \textbf{Project Theme:} & Discussion on the Value and Pressure of Football Defenders \\
    \textbf{Date of Interview:} & Thursday, May 28, 2026 \\
    \textbf{Time of Interview:} & 20:37 -- 20:42 (GMT+08) \\
    \textbf{Interviewer Identifier:} & Speaker 1 Luke Sun\\
    \textbf{Interviewee Identifier:} & Speaker 2 (\redacted) \\
    \textbf{Interviewee Subject Profile:} & School Team Athlete (Central Midfielder)
\end{tabular}

\vspace{2em}
\hrule
\vspace{1.5em}

\section*{Official Verbatim Transcript}

\begin{linenumbers}

\noindent \textbf{Speaker 1 [00:01]:} Okay, you are \redacted, right? Yes. What position do you mainly play on the school football team?

\vspace{0.5em}
\noindent \textbf{Speaker 2 [00:08]:} Central midfielder.

\vspace{0.5em}
\noindent \textbf{Speaker 1 [00:10]:} Central midfielder. How do you define your core responsibilities on the pitch? What do you usually do?

\vspace{0.5em}
\noindent \textbf{Speaker 2 [00:17]:} Participate well in defense, and then orchestrate the team's offense.

\vspace{0.5em}
\noindent \textbf{Speaker 1 [00:21]:} Okay. In the last match, how much of a decisive role do you think you played in the team's victory or defeat? How did you contribute to the victory?

\vspace{0.5em}
\noindent \textbf{Speaker 2 [00:31]:} How much of a decisive role?

\vspace{0.5em}
\noindent \textbf{Speaker 1 [00:33]:} Yes, how did you contribute to the victory? Just...

\vspace{0.5em}
\noindent \textbf{Speaker 2 [00:41]:} Some key passes, I guess.

\vspace{0.5em}
\noindent \textbf{Speaker 1 [00:43]:} Mmh, key passes.

\vspace{0.5em}
\noindent \textbf{Speaker 2 [00:46]:} And playing defense well.

\vspace{0.5em}
\noindent \textbf{Speaker 1 [00:49]:} Do these key passes refer to the attacking side? And defense is defense.

\vspace{0.5em}
\noindent \textbf{Speaker 2 [00:54]:} Yes, yes, yes.

\vspace{0.5em}
\noindent \textbf{Speaker 1 [00:55]:} Okay, if at the beginning of the match, say within the first few minutes, you missed a critical chance, like missing a one-on-one, but you subsequently scored a goal later in the match, do you think people would still mind your earlier mistake?

\vspace{0.5em}
\noindent \textbf{Speaker 2 [01:10]:} I think they wouldn't, right?

\vspace{0.5em}
\noindent \textbf{Speaker 1 [01:11]:} They wouldn't? Is there any reason?

\vspace{0.5em}
\noindent \textbf{Speaker 2 [01:14]:} I feel that since we are a team, everyone makes mistakes at some point. We all need to be tolerant of our teammates.

\vspace{0.5em}
\noindent \textbf{Speaker 1 [01:21]:} Right. But do you think the coach and the spectators would mind the mistake? For instance, the coach, I think they probably would, right? Yes, yes, yes. What about the spectators?

\vspace{0.5em}
\noindent \textbf{Speaker 2 [01:32]:} The spectators, I don't think the spectators would care about that.

\vspace{0.5em}
\noindent \textbf{Speaker 1 [01:34]:} Mmh, okay. When you see a backfield teammate, like \redacted\ or me, making a mistake that leads to conceding a goal, what kind of emotional reaction do you usually have? For instance, we let a ball leak through, and opponents score directly because of it.

\vspace{0.5em}
\noindent \textbf{Speaker 2 [01:50]:} I would feel like, let's just keep defending well, and then try to score a goal back to make up for this mistake.

\vspace{0.5em}
\noindent \textbf{Speaker 1 [01:58]:} Make up for the mistake, meaning using your offense to compensate for some errors on the defensive end. Yes, that works. Yes. Then do you agree with the statement that defensive players face a much higher cost of making mistakes than offensive players?

\vspace{0.5em}
\noindent \textbf{Speaker 2 [02:13]:} I agree with this.

\vspace{0.5em}
\noindent \textbf{Speaker 1 [02:15]:} Is there any reason?

\vspace{0.5em}
\noindent \textbf{Speaker 2 [02:16]:} Because an offensive player's mistake does not directly result in a goal. The probability of it impacting the score directly is very small. But if a defensive player makes a mistake, the probability of it leading to a conceded goal is relatively large.

\vspace{0.5em}
\noindent \textbf{Speaker 1 [02:34]:} Oh, so it's equivalent to saying that their risks are asymmetric. Yes. Then, in post-match summaries or our usual daily discussions, do you discuss goals and assists more frequently, or is it about a defender's passing under pressure, or critical sliding tackles and clearances? Which one has more topicality?

\vspace{0.5em}
\noindent \textbf{Speaker 2 [02:59]:} Has more topicality... I think this depends on the situation on the pitch. For example, if we are completely passive and have been pinned down defending under pressure, we might discuss the defensive issues more. But if we are evenly matched, or if we have a lot of possession in attack, we will discuss the offensive side more.

\vspace{0.5em}
\noindent \textbf{Speaker 1 [03:22]:} What I mean is, in most cases, do goals and assists hold greater topicality for you, or...?

\vspace{0.5em}
\noindent \textbf{Speaker 2 [03:28]:} Then it should be goals and assists.

\vspace{0.5em}
\noindent \textbf{Speaker 1 [03:30]:} Mmh, okay. Do you feel that your own goals or assists bring you more exposure within the campus?

\vspace{0.5em}
\noindent \textbf{Speaker 2 [03:40]:} I think they wouldn't, right?

\vspace{0.5em}
\noindent \textbf{Speaker 1 [03:41]:} They wouldn't? Among fans, do you feel you get talked about more?

\vspace{0.5em}
\noindent \textbf{Speaker 2 [03:49]:} Well, there is still some of that.

\vspace{0.5em}
\noindent \textbf{Speaker 1 [03:51]:} Do you think this extra attention is purely based on the visual attractiveness of the goal, or because a goal itself is inherently more topical?

\vspace{0.5em}
\noindent \textbf{Speaker 2 [04:01]:} I think it's more because a goal itself inherently carries more topicality.

\vspace{0.5em}
\noindent \textbf{Speaker 1 [04:05]:} Mmh, so do you mean that, for example, a sliding tackle is extremely critical, and a goal is also extremely critical, so their actual importance is quite similar, but the goal still edges it out because of its inherent topicality?

\vspace{0.5em}
\noindent \textbf{Speaker 2 [04:20]:} Exactly.

\vspace{0.5em}
\noindent \textbf{Speaker 1 [04:22]:} Do you feel that this kind of inequality, whether it's the inequality of risk or the inequality of topicality, is somewhat unfair to defenders?

\vspace{0.5em}
\noindent \textbf{Speaker 2 [04:37]:} I don't think it's unfair. I used to play as a defender too. I felt that when I played as a defender, my job was simply to do my defensive work well and try not to make mistakes. Now that I have moved to the midfield and frontfield, I focus more of my energy on the offense, hoping to score goals and gain an advantage for the whole team and the club. I think everyone has different responsibilities. What one bears and handles will also be different.

\vspace{0.5em}
\noindent \textbf{Speaker 1 [05:06]:} So you feel that, as of now, defenders are not particularly looked down upon.

\vspace{0.5em}
\noindent \textbf{Speaker 2 [05:10]:} Right, I think everyone belongs to one team, and everyone is equal.

\vspace{0.5em}
\noindent \textbf{Speaker 1 [05:13]:} Mmh, okay. Then this concludes it. OK. OK, thank you.

\end{linenumbers}

\end{document}